\section{Mosaic, version one}
\label{sec:ch02-mosaic-v1}
\index{Mosaic|(}

Mosaic sells furniture, lighting, kitchen equipment, textiles and storage. The
seeded catalog has twenty-five products in it, which is enough to make a list
query interesting and few enough that you can read the whole thing in one
listing if you want to.

Catalog owns what a product \emph{is}:

\begin{minted}{csharp}
public sealed record Product(
    Guid Id,
    string Sku,
    string Title,
    string? Description,
    ProductCategory Category);
\end{minted}

Five properties, and none of them is the price. That omission is the whole
design and I will come back to it in the next section.

\index{Mosaic!in-memory data}%
The data lives in memory, in a singleton holding a list built from a C\#
collection expression. There is no database in this chapter.
Chapter~\ref{ch:data-without-the-n-plus-1} brings EF~Core and Postgres, and
doing it here would mean spending a third of a chapter called \emph{quickly} on
connection strings. Product identifiers come from a fixed pattern rather than
being random:

\begin{minted}{csharp}
public static Guid ProductId(int n) =>
    new($"a0000000-0000-4000-8000-{n:D12}");
\end{minted}

That looks like a shortcut and it is doing real work. The other five domains
need to refer to products, and each of them writes its own copy of this helper
rather than calling Catalog's. It is duplication, deliberately, because after
chapter~\ref{ch:strangling-the-monolith} these domains are separate services
that cannot share a file, and the identifier is the only contract left between
them. Seed data that models the eventual boundary is easier to split than seed
data that models a shared C\# assembly.

\index{Mosaic!CatalogService}%
Between the data and the graph sits a service:

\begin{minted}{csharp}
public sealed class CatalogService(
    InMemoryCatalogData data,
    ServiceCallCounter counter)
{
    public async Task<Product?> GetProductByIdAsync(
        Guid id,
        CancellationToken cancellationToken)
    {
        await counter.RecordLookupAsync(cancellationToken);
        return data.Products.FirstOrDefault(p => p.Id == id);
    }
}
\end{minted}

Three things about that method are choices rather than accidents. It takes one
key and answers about one thing, and there is no overload anywhere in Mosaic
that accepts a list of identifiers. It returns a \code{Task} even though a
linear scan of a list in memory is synchronous, so that the signature does not
have to change when the work becomes genuinely asynchronous. And it tells a
counter that a lookup happened, which is how
section~\ref{sec:ch02-still-wrong} gets its number.

The root fields are a static class with the \code{[QueryType]} attribute:

\begin{minted}{csharp}
[QueryType]
public static partial class CatalogQueries
{
    public static Task<IReadOnlyList<Product>> GetProductsAsync(
        CatalogService catalog,
        CancellationToken cancellationToken)
        => catalog.GetProductsAsync(cancellationToken);

    public static Task<Product?> GetProductByIdAsync(
        [ID] Guid id,
        CatalogService catalog,
        CancellationToken cancellationToken)
        => catalog.GetProductByIdAsync(id, cancellationToken);
}
\end{minted}

\index{HotChocolate!resolver parameters}%
Note what the parameters are doing. \code{id} is a GraphQL argument,
\code{catalog} is a service resolved from the container, and
\code{cancellationToken} is neither. HotChocolate decides which is which by
looking at the type and at what is registered in dependency injection, and it
gets this right almost always. When it gets it wrong the failure is memorable:
forget to register a service and HotChocolate concludes the parameter must be an
argument, tries to build an input type out of it, and the schema fails to build
with \enquote{No compatible constructor found for input type}. The service is
not missing, in its view. It has been promoted.

\index{HotChocolate!ID scalar}%
The last piece is a single-field type extension:

\begin{minted}{csharp}
[ObjectType<Product>]
public static partial class ProductNode
{
    [ID]
    public static Guid GetId([Parent] Product product) => product.Id;
}
\end{minted}

That is the same \code{[ID]} correction from
section~\ref{sec:ch02-two-ways}, and the new part is \code{[Parent]}, which
marks the object the field is being resolved on. It is a small thing, but it is
the first example of the pattern that the next section is built on: a class
whose entire job is to add one field to a type declared somewhere else.

Ask for a product and you get one:

\begin{minted}{graphql}
{
  productById(id: "a0000000-0000-4000-8000-000000000001") {
    sku
    title
    category
  }
}
\end{minted}

\begin{minted}{json}
{
  "data": {
    "productById": {
      "sku": "MOS-FRN-0001",
      "title": "Larsen Oak Dining Table",
      "category": "FURNITURE"
    }
  }
}
\end{minted}

That is Catalog, complete. It is also, on its own, an entirely ordinary GraphQL
service of the kind chapter~\ref{ch:why-one-graph-is-never-enough} said was the
right first move. The interesting part starts when the other five domains want a
say in what a product looks like.
